\section{幻灯片图谱与关键结构}

\subsection{策略角色分布}

\begin{figure}[H]
\centering
\includegraphics[width=0.92\textwidth,page=3]{08_cs224r_reward_learning_2025.pdf}
\caption{策略与 reward score 的流程图（slide 3）：展示 reward model、policy、evaluators 之间的交互。}
\end{figure}
\newpage

\subsection{数据流与日志}

\begin{figure}[H]
\centering
\includegraphics[width=0.92\textwidth,page=4]{08_cs224r_reward_learning_2025.pdf}
\caption{数据与日志流：从 human label 到 logging dashboard 的多层管道。}
\end{figure}
\newpage

\subsection{Policy Rollout}

\begin{figure}[H]
\centering
\includegraphics[width=0.92\textwidth,page=5]{08_cs224r_reward_learning_2025.pdf}
\caption{Policy rollout 与 reward trace 交互图，是之前章节中可观测性 checklist 的直观延伸。}
\end{figure}
\newpage

\subsection{KL 罚项与温度调度}

\begin{figure}[H]
\centering
\includegraphics[width=0.92\textwidth,page=6]{08_cs224r_reward_learning_2025.pdf}
\caption{KL penalty 及 temperature schedule 曲线示意，用于支持 KL penalty 的解释章节。}
\end{figure}
\newpage

\subsection{Fail-safe 监控}

\begin{figure}[H]
\centering
\includegraphics[width=0.92\textwidth,page=8]{08_cs224r_reward_learning_2025.pdf}
\caption{Failure monitor dashboard 例图，对应本章的报警与 rollout check list。}
\end{figure}
\newpage

\subsection{人类偏好比较流程}

\begin{figure}[H]
\centering
\includegraphics[width=0.92\textwidth,page=9]{08_cs224r_reward_learning_2025.pdf}
\caption{Preference comparison pipeline，highlight human-in-the-loop steps。}
\end{figure}
\newpage

\subsection{Sim-to-Real 图示}

\begin{figure}[H]
\centering
\includegraphics[width=0.92\textwidth,page=10]{08_cs224r_reward_learning_2025.pdf}
\caption{Sim-to-Real checklist：强调 domain randomization 和 reward smoothing。}
\end{figure}

\subsection{本章小结}

Slide gallery 让我们看到 reward learning 将训练、评估、监控串成一体的全景图，方便对照文中 checklist 和操作流水线。

